\begin{frame}{Dyna-Q 복습: 학습된 모델로 가상 경험 만들기}
  Week 7.3의 \kb{Dyna-Q}는 세 가지를 \textbf{동시에} 했다:
  \begin{enumerate}
    \item 실제 환경과 상호작용해 \textbf{경험}을 쌓는다
    \item 그 경험으로 간단한 모델 $\hat{P}(s'|s,a)$, $\hat{R}(s,a)$를
      \textbf{학습}한다
    \item 그 모델로 \textbf{가상의 추가 경험}을 만들어 Q값을
      더 많이 갱신한다
  \end{enumerate}
  \vskip0.6em
  \begin{block}{모델기반 강화학습(Model-Based RL)}
    Dyna-Q 아이디어의 일반화 --- 환경의 모델(전이확률 + 보상함수,
    또는 그 근사)을 학습하거나 이미 알고 있다면, ``실제로 행동해보지
    않고도'' 여러 수를 미리 시뮬레이션할 수 있다.
  \end{block}
  \vskip0.5em
  \begin{itemize}
    \item Dyna-Q는 이 시뮬레이션을 \textbf{값 갱신}에 썼다
    \item 이번 장의 MCTS는 \textbf{지금 당장 어떤 행동을 할지 결정하는
      데} 직접 쓴다 --- 이 점이 차이
  \end{itemize}
\end{frame}
